\documentclass[../../../include/open-logic-section]{subfiles}

\begin{document}

\olfileid[fr]{sth}{cardinals}{cp}
\olsection{Principe de Cantor}

Revenons à la \olref[ordinals][vn]{sec}. Nous parlions
des ensembles bien ordonnés et suggérions qu'il serait
utile de disposer d'objets tenant lieu de représentants
des bons ordres. Dans cette perspective, nous avons
introduit les ordinaux, puis montré dans le
\olref[ordinals][ordtype]{ordtypesworklikeyouwant}
qu'ils se comportent comme nous le souhaitions :
\[
	\ordtype{A, <} = \ordtype{B, \lessdot}
	\text{ ssi } \tuple{A, <}
	\isomorphic \tuple{B, \lessdot}.
\]
Revenons encore plus loin, à la
\olref[sfr][siz][equ]{sec}. Nous y avons
introduit, de manière naïve, la notion de
«~taille~» d'un ensemble. Plus précisément,
nous avons dit que deux ensembles sont équipotents,
$\cardeq{A}{B}$, s'il existe !!a{bijection}
$f \colon A \to B$. Cette notion est
intrinsèquement {plus simple} que celle de bon
ordre : seules les !!{bijection}s nous intéressent,
et non, comme pour les isomorphismes d'ordres,
la question de savoir si ces !!{bijection}s
«~préservent une structure~».

Une idée naturelle se présente alors. De même
que nous avons introduit certains objets, les
\emph{ordinaux}, pour mesurer les bons ordres,
nous pouvons introduire d'autres objets, les
\emph{cardinaux}, pour mesurer la taille.
Tel est le but de ce chapitre.

Avant de dire ce que seront ces cardinaux,
énonçons un principe qu'ils devraient satisfaire.
Si nous écrivons $\card{X}$ pour la cardinalité
de l'ensemble $X$, nous voudrions avoir :
\[
	\card{A} = \card{B}
	\text{ ssi } \cardeq{A}{B}.
\]
Nous l'appellerons le \emph{principe de Cantor},
car Cantor fut probablement le premier à l'avoir
très clairement à l'esprit. (Nous reviendrons
sur sa relation avec le \emph{principe de Hume}
dans la \olref[hp]{sec}.) Notre objectif est donc
de définir $\card{X}$ pour chaque $X$ de façon
à satisfaire le principe de Cantor.

\end{document}
